\section{A checked critical-line zeta bound}

The following input is obtained by specializing the checked logarithmic
phase estimates imported from RiemannGaussian.  It concerns the actual
Mathlib zeta function, and assumes no growth estimate for that function.
The imported proof slice and its exact upstream revision are recorded in
\texttt{docs/rg-extraction-manifest.json}.

\begin{theorem}[Weyl growth on the critical line]
\label{construction:weyl-zeta}
\lean{Sidon.riemannZeta_critical_line_weyl,Sidon.riemannZeta_critical_line_power,Sidon.exists_riemannZeta_critical_growth}\leanok
For $s=\tfrac12+it$ and $|t|\geq2$,
\[
 |\zeta(s)|\leq8192|t|^{1/6}\log|t|
 \leq393216|t|^{3/16}.
\]
There is a constant $C>0$ such that for every real $t$,
\[
 |\zeta(\tfrac12+it)|\leq C(1+|t|)^{3/16}.
\]
\end{theorem}
\begin{proof}
The imported generic theorem uses
\[
 \alpha_k=\frac1{2^{k+2}-2},\qquad
 \sigma_k=1-(k+2)\alpha_k,
\]
and proves $|\zeta(\sigma_k+it)|\leq8192|t|^{\alpha_k}\log|t|$
for every integer $k\geq1$ and $|t|\geq2$.
At $k=1$ these definitions give $\alpha_1=1/6$ and $\sigma_1=1/2$,
so its hypotheses are exactly the hypotheses of the first estimate.

To specify the analytic content of this input, its Euler truncation
chooses a dyadic integer $M$ with $|s|<M\leq5t$ when $t\geq2$.
The checked uniform Euler remainder bounds zeta by the norm of the
positive Dirichlet prefix through $M-1$, plus $6$.
The prefix is decomposed into complete dyadic blocks.  Each block is
bounded by the checked derivative estimates for the actual phase
$-t\log n$, and their finite budget is at most a constant times
$t^{\alpha_k}\log t$ on the line $\sigma_k$.
The generic theorem combines these two estimates; $t^{\alpha_k}\geq1$
and $\log t\geq\log2>1/2$ absorb the additive $6$ in the stated
constant $8192$.  For negative $t$, apply the positive-height result to
$\overline{s}$ and use
$\zeta(\overline{s})=\overline{\zeta(s)}$ and equality of the norms.
This is a finite Euler truncation argument; it does not assume an
approximate functional equation or a critical-line growth hypothesis.

For the second estimate, the elementary inequality
$\log x\leq x^{b-a}/(b-a)$, valid for $x>0$ and $a<b$, gives
\[
 x^{1/6}\log x\leq48x^{3/16},
 \qquad \frac3{16}-\frac16=\frac1{48}.
\]
Multiplying by $8192$ gives $393216$.
The point $1/2+it$ never equals the pole $1$; Mathlib's differentiability
of zeta away from its pole therefore proves continuity on this line.
Compactness of $[-2,2]$ supplies a constant $B\geq1$ bounding its norm
there.  Take $C=\max(B,393216)>0$.
For $|t|\leq2$, use $(1+|t|)^{3/16}\geq1$.
For $|t|\geq2$, use the tail estimate and monotonicity of the positive
real power to replace $|t|^{3/16}$ by $(1+|t|)^{3/16}$.
These cases cover every real $t$, including their common endpoints.
\end{proof}
